\documentclass[11pt]{article}
% Standalone preamble for web-only supplementary notes.
\usepackage[T1]{fontenc}
\usepackage{amsmath,amssymb,amsthm}
\usepackage{hyperref}
\usepackage{cleveref}
\DeclareMathOperator{\Map}{Map}
\DeclareMathOperator{\Fun}{Fun}
\newcommand{\iso}{\xrightarrow{\sim}}
\newtheorem{lemma}{Supplementary Lemma}
\newtheorem{corollary}[lemma]{Supplementary Corollary}
\crefname{lemma}{Supplementary Lemma}{Supplementary Lemmas}
\Crefname{lemma}{Supplementary Lemma}{Supplementary Lemmas}
\crefname{corollary}{Supplementary Corollary}{Supplementary Corollaries}
\Crefname{corollary}{Supplementary Corollary}{Supplementary Corollaries}

\begin{document}
\section*{Decoding mapping animae and restriction}

\href{mapping-animae.html\#con:Functors_As_Points_Of_Mapping_Animae}{Naming and decoding}
relate functors $C\to D$ to absolute objects of $\Map(C,D)$.
When this correspondence is used in a universal-property argument, one
also needs its action on identification animae and its compatibility with
restriction. This note explains those two ingredients.

Put $j_C=(!,\mathrm{id}_C)\colon C\to\star\times C$. For an absolute
object $u\colon\star\to\Map(C,D)$, write
\[
 D(u)=\operatorname{uncurry}(u)\circ j_C\colon C\to D.
\]
Here $D(u)$ abbreviates the decoded functor in this note. For
$i\colon B\to C$, write $i^*\colon\Map(C,D)\to\Map(B,D)$ for
precomposition. All the categories and endpoints in the following
statements are absolute.

\begin{lemma}
\label[lemma]{lem:Supplementary_Decoding_Equivalence}
For absolute objects $u,v$ of $\Map(C,D)$, decoding induces an equivalence
\[
 D_{u,v}\colon(u=v)\longrightarrow(D(u)=D(v)).
\]
Consequently, each specified identification $\delta\colon D(u)=D(v)$
admits an identification $\widetilde\delta\colon u=v$ equipped with
a comparison $D_{u,v}(\widetilde\delta)=\delta$.
\end{lemma}
\begin{proof}[Proof idea]
The functor $D_{u,v}$ is the composite of the action of uncurrying on
identifications with prewhiskering by $j_C$. The first is an equivalence
by the mapping-anima axiom, applied to the anima $\star$.
The second is an equivalence because $j_C$ is an equivalence: its inverse
is the second projection, with the two comparisons obtained from the
product and terminal universal properties. Whiskering by an equivalence
is treated in \href{supplementary-changing-endpoints.html}{Changing endpoints
in identification animae}. Applying a chosen inverse of $D_{u,v}$ gives
the lift and its specified image comparison.
\end{proof}

\begin{lemma}
\label[lemma]{lem:Supplementary_Decoding_Restriction}
For $i\colon B\to C$ and an absolute object $u$ of $\Map(C,D)$,
there is a specified identification
\[
 e_i(u)\colon D(i^*\circ u)=D(u)\circ i.
\]
\end{lemma}
\begin{proof}[Proof idea]
First compare the two functors from $B$ to $\star\times C$:
\[
 (\mathrm{id}_{\star}\times i)\circ j_B=j_C\circ i.
\]
Their first projections are identified by the terminal universal
property. Their second projections are identified with $i$ by the
projection comparisons and unitors. The product universal property
therefore supplies the displayed identification.

Uncurrying $i^*\circ u$ gives
$\operatorname{uncurry}(u)\circ(\mathrm{id}_{\star}\times i)$.
Precompose this comparison with $j_B$, reassociate, apply the preceding
product comparison, and reassociate back. The resulting chain is the
chosen $e_i(u)$.
\end{proof}

\section*{The square on identification animae}

Write $\star$ for whiskering and $\cdot$ for vertical composition,
with the rightmost factor applied first. A bar denotes a constant family.

\begin{lemma}
\label[lemma]{lem:Supplementary_Decoding_Naturality}
For $i\colon B\to C$, absolute objects $u,v$ of $\Map(C,D)$,
and a functor $\gamma\colon A\to(u=v)$, one has an identification
of functors on $A$
\[
\begin{aligned}
 &\overline{e_i(v)}\cdot
       \bigl(D_{i^*u,i^*v}\circ(i^*\star\gamma)\bigr)\\
 &\qquad=\bigl((D_{u,v}\circ\gamma)\star i\bigr)
                    \cdot\overline{e_i(u)}.
\end{aligned}
\]
Here $i^*u$ abbreviates $i^*\circ u$. Both sides take values in
$(D(i^*u)=D(v)\circ i)$.
\end{lemma}
\begin{proof}[Proof idea]
Follow the four stages constructing $e_i$ in
\Cref{lem:Supplementary_Decoding_Restriction}, now with the family
$\gamma$ inserted. Compatibility of uncurrying with precomposition gives
the first square. Compatibility of prewhiskering with composition gives
the reassociation square. Fixed-inner interchange gives the square for
the comparison $(\mathrm{id}_{\star}\times i)\circ j_B=j_C\circ i$.
A final reassociation square handles the inverse associator at the end.

Paste the four squares. This first proves the formula for the explicit
family obtained by uncurrying $\gamma$ and then prewhiskering by $j_C$.
The evaluation comparison for the composite defining $D_{u,v}$ identifies
that family with $D_{u,v}\circ\gamma$, yielding the displayed formula.
\end{proof}

Taking $A=(u=v)$ and $\gamma=\mathrm{id}_{(u=v)}$, and removing the
resulting identity precompositions by the unit comparisons, gives the square
between the actual restriction functors on identification animae.
This is stronger than comparing their values on individual absolute
identifications: it is the square used to transfer equivalence of
restriction from mapping animae to decoded functors. Its application to
coproducts is explained in
\href{supplementary-identifications-on-coproducts.html}{Identifications on coproducts}.
\end{document}
